% Add bibliography key BlasiakRep for Jonah Blasiak,
% Representation theory of the nonstandard Hecke algebra,
% arXiv:1201.2209v2 (2012), https://arxiv.org/abs/1201.2209v2.
\subsection{An actual nonstandard tower: specialization fails on lines and succeeds after mixing}
There is already a nonstandard tower to which the generic path mechanism
applies. Blasiak \cite[Theorem~5.1, p.~14; Proposition~5.2, pp.~14--15;
Definition~6.1, p.~23]{BlasiakRep} proves that the two-row nonstandard
Temperley--Lieb algebra over $K=\Q(v)$ is split semisimple, that its
successive restrictions are multiplicity free, and hence that its simple
modules have seminormal path lines. These are established source results,
not new consequences of purity. The source distinguishes these lines
from a basis compatible with classical decomposition at $v=1$; its
Remark~6.2 also gives a non-multiplicity-free example outside the
two-row quotient. In particular, its simple module at \emph{tower rank
four}
\[
 M_K=M_{(3,1),K}\otimes M_{(2,2),K}
\]
has dimension six and restriction
\begin{equation}\label{eq:towerrestriction}
 M_K|_{A_{3,K}}=N(\delta)\oplus N(D-2)\oplus\epsilon_+
                  \oplus\epsilon_-.
\end{equation}
Here the rank-three notation is that of \fid{F19}/\fid{F21}.
The following comparison computes exactly what happens to its path lines
in the natural ordinary Specht tensor lattice.

\begin{finding}[\fid{F30}: regular path mixing and a projector base-change defect]\label{f:tower}
Let $R=\Q[v,v^{-1}]_{(v-1)}$, $k=R/(v-1)=\Q$, and fix
\[
 M_R=\mathsf S_{(3,1),R}\otimes_R\mathsf S_{(2,2),R},
\]
with the specified nonstandard action. Write $A_R$ for the actual
nonstandard algebra acting at tower rank four. Its specialization has
\[
 M_1\simeq[(3,1)]\oplus[(2,1,1)].
\]
The two generic path lines through the penultimate summand $N(\delta)$
lie in a regular saturated rank-two plane $E_R$. No vector on either
line which is diagonally rescaled to be regular with nonzero reduction
can specialize into one of these two classical constituents.

Nevertheless a regular non-diagonal change of the six path representatives
in this same lattice does produce a specialization-compatible basis.
It mixes only two pairs of corresponding rank-two path labels, by
unit rescaling followed by the blocks
\[
 \begin{pmatrix}1&1\\1&-1\end{pmatrix}.
\]
Finally
\[
 \End_{A_R}(M_R)=R,\qquad \End_{A_1}(M_1)=k\times k,
\]
and the natural base-change map is $a\mapsto(a,a)$. Thus the two
classical projectors do not lift to endomorphisms commuting with the
whole nonstandard deformation, although the stated regular mixed basis
exists.
\end{finding}
\begin{proof}
\par\smallskip\noindent\textbf{A fixed saturated plane.}\quad
The ordinary parabolic Hecke algebra $H_3\subset H_4$ has trivial
projector
\[
 p_+=\frac{\sum_{w\in S_3}v^{\ell(w)}T_w}
              {\sum_{w\in S_3}v^{2\ell(w)}}.
\]
The denominator specializes to six and is a unit of $R$. The numerator
is multiplied by $v$ by each simple generator on either side; squaring
it therefore multiplies it by the displayed denominator. Hence $p_+$
is a regular central idempotent of $H_3$. Ordinary branching gives
\[
 M_{(3,1)}|_{H_3}=M_{(3)}\oplus M_{(2,1)},\qquad
 M_{(2,2)}|_{H_3}=M_{(2,1)}.
\]
It follows that
\[
 E_R=M_R(p_+\otimes1)
\]
is a free rank-two direct summand, stable under $A_{3,R}$, and its
generic fiber is exactly the $N(\delta)$ summand in
\eqref{eq:towerrestriction}. Because the projector is regular,
\begin{equation}\label{eq:towersaturated}
 E_R=E_K\cap M_R.
\end{equation}
Indeed an element of the intersection is fixed by $p_+\otimes1$
and so lies in its $R$-image. Every regular nonzero reduction of a
rescaled generic path vector in $E_K$ therefore lies in $E_1$.

\par\smallskip\noindent\textbf{The special decomposition and its class sum.}\quad
At $v=1$, $P_i=2(1+s_i)$ for the diagonal symmetric-group action,
so the specialized algebra image generates that action. On the five
cycle types $(1^4),(2,1,1),(2,2),(3,1),(4)$, the characters are
\[
 \chi_{31}=(3,1,-1,0,-1),\qquad
 \chi_{22}=(2,0,2,-1,0),\qquad
 \chi_{211}=(3,-1,-1,0,1).
\]
Their product satisfies
\[
 \chi_{31}\chi_{22}=(6,0,-2,0,0)=\chi_{31}+\chi_{211},
\]
which proves the stated classical decomposition in characteristic zero.
Let $U$ be the sum-zero subspace of $\mathbb R^4$, and $W$ the
sum-zero subspace of the permutation representation on the three
pairings of four letters. These are orthogonal models for $[31]$
and $[22]$. The $S_3$ fixing the fourth letter has invariant line
in $U$ generated by
\[
 w=(1,1,1,-3),\qquad \langle w,w\rangle=12.
\]
Thus $E_1=\mathbb Rw\otimes W$. Let $p$ project orthogonally onto
$\mathbb Rw$, put $\Pi=p\otimes1$, and let
\[
 Z=\sum_{1\le i<j\le4}(ij)
\]
act diagonally. For $i,j\le3$, the transposition fixes $w$ and
$p\rho_U((ij))p=p$. For an external transposition $(i4)$ its inner
product with $w$ is $-4$, giving
\[
 p\rho_U((i4))p=-\tfrac13p.
\]
The action on the three pairings is the quotient $S_4\to S_3$ with
Klein-four kernel. The three internal transpositions map bijectively
to its three transpositions, and so do the three external ones.
The sum of those three matrices in the natural three-point
permutation module is the all-ones matrix; it vanishes on $W$.
Consequently both transposition sums on $W$ are zero and
\begin{equation}\label{eq:towercompression}
 \Pi Z\Pi=p\otimes(0-\tfrac13\,0)=0.
\end{equation}

The class sum $Z$ acts on $[31]$ and $[211]$ by $+2$ and $-2$:
its scalar is six times the transposition character divided by the
dimension, namely $6(\pm1)/3$. If a nonzero $x\in E_1$ belonged
to either constituent, then $\Pi Z\Pi x=\pm2x$, contradicting
\eqref{eq:towercompression}. Thus $E_1$ meets each constituent only
in zero. Together with \eqref{eq:towersaturated}, this proves the
diagonal-normalization obstruction. All identities in this real
model have rational coefficients, so the conclusion holds over $k$.

\par\smallskip\noindent\textbf{Regular complementary path planes.}\quad
Put $J_R=M_R((1-p_+)\otimes1)$. It is the complementary regular
rank-four summand. The three central block idempotents
$z_+,z_-,z_{\rm mat}$ of the actual rank-three algebra in \fid{F19}
are defined over $R$, without dividing by the colliding-root
difference. They commute with the ordinary parabolic projector.
Set
\[
 L_{+,R}=J_Rz_+,\quad L_{-,R}=J_Rz_-,\quad F_R=J_Rz_{\rm mat}.
\]
These direct summands are free over the local ring $R$, with generic
types $\epsilon_+,\epsilon_-,N(D-2)$ and ranks $1,1,2$. Hence
\[
 M_R=L_{+,R}\oplus L_{-,R}\oplus E_R\oplus F_R.
\]
This separates $E_R$ from $J_R$ by the ordinary projector first;
it does not try to separate the two colliding matrix branches by a
singular nonstandard central projector.

At specialization the four $S_3$ types are respectively trivial,
sign, standard, standard. The first-factor orthogonal complement
$U_0=w^\perp\subset U$ is standard, so
$J_1=U_0\otimes W=1\oplus\mathrm{sgn}\oplus[21]$ and $F_1$
is its unique standard summand. Write these specialized spaces as
$L_+,L_-,E,F$. Since $Z$ commutes with $S_3$, it sends $E$ into
$E\oplus F$; there is no homomorphism from the standard module
to either scalar type. Equation~\eqref{eq:towercompression} removes
the $E$ component. Also $Z^2=4$ by the classical decomposition.
It follows that
\[
 Z(E)=F,\qquad Z(F)=E,\qquad J:=Z/2:E\longrightarrow F
\]
is an isomorphism with inverse $Z/2$. Its graph planes
\[
 G_+=\{x+Jx:x\in E\},\qquad G_-=\{x-Jx:x\in E\}
\]
have $Z$-eigenvalues $+2$ and $-2$. The classical restrictions
$[31]|_{S_3}=1\oplus[21]$ and
$[211]|_{S_3}=\mathrm{sgn}\oplus[21]$ therefore identify the
actual constituents as
\begin{equation}\label{eq:towergraphs}
 [31]=L_+\oplus G_+,\qquad [211]=L_-\oplus G_-.
\end{equation}

\par\smallskip\noindent\textbf{A two-pair mixing repair.}\quad
The idempotent $e=P_1/D$ from tower rank two is regular. Each of $E_R,F_R$
splits by $e,1-e$ into two rank-one free direct summands. Choose
bases $x_{+,R},x_{-,R}$ of these two lines in $E_R$, and
$f_{+,R},f_{-,R}$ in $F_R$. Their generic lines are exactly the
four tower paths through the two standard penultimate labels.
At specialization $e=(1+s_1)/2$. Since $J$ commutes with $s_1$,
it carries each line of $E$ to the matching line of $F$. Rescale
each $f_{\sigma,R}$ by a constant unit of $\Q$ so that its
reduction is $Jx_\sigma$.
For each $\sigma=+,-$, replace the two regular vectors by
\[
 g_{\sigma,+}=x_{\sigma,R}+f_{\sigma,R},\qquad
 g_{\sigma,-}=x_{\sigma,R}-f_{\sigma,R}.
\]
The displayed two-by-two block has determinant $-2$, a unit in
$R$, so these four vectors, together with bases of $L_{+,R}$
and $L_{-,R}$, form an $R$-basis. Their reductions lie in the
two graph planes in \eqref{eq:towergraphs}. This proves the
regular specialization-compatible basis claim in the original
fixed lattice. It uses non-diagonal mixing across distinct generic
penultimate summands, precisely the freedom left open by the first
part of the finding.

\par\smallskip\noindent\textbf{Why the classical projectors still do not lift.}\quad
By the cited split simplicity,
$\End_{A_K}(M_K)=K\,\id$. Any element of $\End_{A_R}(M_R)$
therefore becomes a scalar $\lambda\id$ generically. In an
$R$-basis its diagonal entries belong to $R$, forcing $\lambda\in R$.
Conversely all such scalars commute with the action. This proves
$\End_{A_R}(M_R)=R$. At specialization the image is the diagonal
$k[S_4]$ action with two pairwise nonisomorphic absolutely simple
constituents, each once, so its endomorphism ring is $k\times k$.
The base-change map sends a scalar to $(a,a)$; it is injective with
one-dimensional cokernel, and neither nontrivial classical projector
lies in its image.

The three-vector groups in the mixed basis span direct summands of
the underlying free $R$-module, whose reductions are the two
classical constituents. They are not $A_R$-stable generically,
which would contradict generic simplicity. Thus the constructive
basis comparison makes no false claim of a nonstandard-equivariant
projector lift.
\end{proof}

\paragraph{Exact scope and the coherent dictionary.}
The excluded operation is diagonal normalization of this tower's
path lines in the specified ordinary Specht tensor lattice; the
positive operation is the explicit regular mixing just constructed.
No arbitrary-lattice obstruction, VIII canonical-basis identification,
uniform general specialization algorithm or positive structure-constant
rule is inferred. The known path labels count generic simple dimensions,
not classical Kronecker multiplicities by themselves. Their verification
is polynomial in the expanded tower rank, not automatically in its
binary logarithm.

In the language of schemes, $M_R$ is a vector bundle on $\Spec R$
with an action of a finite coherent algebra. Images of the regular
idempotents above are coherent direct summands and hence vector
subbundles. The regular mixing is an automorphism of the underlying
bundle; it does not commute with the whole algebra action. The coherent
commutant is the kernel of the finitely many linear commutation
equations for the generators and is the scalar sheaf in this case.
Its formation does not commute with the special fiber: that fiber
has a second classical projector. This is an explicit base-change
phenomenon for coherent endomorphisms, compatible with Grothendieck's
scheme and module language. It supplies neither a Frobenius action nor
a weight statement, and purity alone does not decide which algebra
equivariance a proposed comparison should retain.
